% 
% ======================================================================
\RequirePackage{docswitch}
% \flag is set by the user, through the makefile:
%    make note
%    make apj
% etc.
\setjournal{\flag}

\documentclass[\docopts]{\docclass}

% You could also define the document class directly
%\documentclass[]{emulateapj}

% Custom commands from LSST DESC, see texmf/styles/lsstdesc_macros.sty
\usepackage{lsstdesc_macros}

\usepackage{graphicx}
\graphicspath{{./}{./figures/}}
\bibliographystyle{apj}

% Add your own macros here:

\newcommand{\textul}{\underline}

% 
% ======================================================================

\begin{document}

\title{ Approximating photo-z PDFs for large surveys }

\maketitlepre

\begin{abstract}

Upcoming and ongoing galaxy surveys will produce redshift probability distribution functions (PDFs) in addition to traditional photometric redshift (photo-$z$) point estimates.  However, the storage of photo-$z$ PDFs may present a challenge as the dataset size increases, as we face a trade-off between the accuracy of subsequent science measurements and the storage cost. We investigate a number of different PDF approximations, using metrics that quantify performance in both large scale structure and weak gravitational lensing studies. In the process, we present \texttt{qp}, a Python library enabling the evaluation of various approximations of 1-dimensional PDFs, as suitable for photometric redshifts.

\end{abstract}

% Keywords are ignored in the LSST DESC Note style:
\dockeys{latex: templates, papers: awesome}

\maketitlepost

% ----------------------------------------------------------------------
% 

%\begin{figure}
%\includegraphics[width=0.75\columnwidth]{outline.png}
%\caption{}. \label{fig:outline}}
%\end{figure}

\section{Introduction}
\label{sec:intro}

%This is a paper and note template for the LSST DESC \citep{Overview,ScienceBook,WhitePaper}.
%You can delete all this tutorial text whenever you like.
%
%You can easily switch between various \LaTeX\xspace styles for internal notes and peer reviewed journals.
%Documents can be compiled using the provided \code{Makefile}.
%The command \code{make} with no arguments compiles \code{main.tex} using the  \code{lsstdescnote.cls} style.
%If you want to upgrade your Note into a journal article, just choose a journal name, between \code{make apj} (ApJ preprint format), \code{make apjl} (which uses the \code{emulateapj} style), \code{make prd}, \code{make prl}, and \code{make mnras}.

[Problem statement: photo-zs, LSST]


% ----------------------------------------------------------------------

%\section{Commands}
%\label{sec:commands}

%There are a number of useful \LaTeX\xspace commands predefined in \code{macros.tex}.
%Notice that the section labels are prefixed with \code{sec:} to allow the use of the \verb=\secref= command to reference a section (\ie, \secref{intro}).
%Figures can be referenced with the \verb=\figref= command, which assumes that the figure label is prefixed with \code{fig:}.
%In \figref{example} we show an example figure.
%You'll notice that the actual figure file is found in the \code{figures} directory.
%However, because we have specified this directory in our \verb=\graphicspath= we do not need to explicitly specify the path to the image.
%
%The \code{macros.tex} package also contains some conventional scientific units like \angstrom, \GeV, \Msun, etc. and some editorial tools for highlighting \FIXME{issues}, \CHECK{text to be checked}, \COMMENT{comments}, and \NEW{new additions}.


% ----------------------------------------------------------------------

\section{Methods}
\label{sec:methods}

[Describe possible interpolation schemes used in conversions for both approximations and metrics]

%Similar to the figure before, here we have included a table of data from \code{tables/table.tex}.
%Notice that again we are able to reference \tabref{example} with the \verb=\tabref= command using the \code{tab:} prefix.
%Also notice that we haven't needed to specify the full path to the table because in the \code{Makefile} we include \code{./tables} directory in the \code{\$TEXINPUTS} environment variable.
%
%\begin{table}
  \begin{center}
  \caption{Example table. \label{tab:example}}
  %\begin{ruledtabular}
  \begin{tabular}{lccc}
\hline\hline
Column 1 & Column 2 & Column 3 &  Column 4 \\[3pt]  
     &    $\deg$     & $\kpc$   &  $\deg$ \\[4pt]
\hline
Obj1 & (0,0) & 10 & 0.1 \\
... & ... & ... & ... \\
ObjN & (0,0) & 10 & 0.1
\\\hline\hline
\end{tabular}
\end{center}
%\end{ruledtabular}
\end{table}

%\begin{\tabletype}{l ccccccc }
%\tablewidth{0pt}
%\tabletypesize{\tiny}
%\tablecaption{ An example table. \label{tab:example}}
%\tablehead{
%(1) & (2) & (3) & (4) & (5) & (6) & (7) & (8)\\
%Name & GLON,GLAT & Distance & $r_{1/2}$ & $\log_{10}(J_{\rm meas})$ & $\log_{10}(J_{\rm pred})$ & Sample & Refrence \\
% & (deg) & (kpc) & (pc) & $\log_{10}(\GeV^2 \cm^{-5})$ & $\log_{10}(\GeV^2 \cm^{-5})$ & & 
%}
%\startdata
%Bootes I                     & 358.08,69.62   & 66  & 189  & $18.8 \pm 0.2$ & 18.5           & I,S,C & ... \\
%\\
%...\\
%\\
%Willman 1                    & 158.58,56.78   & 38  & 19   & $19.1 \pm 0.3$ & 18.9           & I,S & ... \\
%\enddata
%{\footnotesize \tablecomments{ (1) The first column. (2) The second column ...}}
%\end{\tabletype}

%
%Equations appear as follows, and can be referred to as, for example, \eqnref{example} -- just as for tables, we use the \verb=\eqnref= command using the \code{eqn:} prefix.
%\begin{equation}
%  \label{eqn:example}
%  \langle f(k) \rangle = \frac{ \sum_{t=0}^{N}f(t,k) }{N}
%\end{equation}

\subsection{Approximation Methods}
\label{sec:approx}

[Lay out all parametrizations, pros and cons]

\subsubsection{Samples}
\label{sec:samples}

\subsubsection{Regular Binning}
\label{sec:bins}

\subsection{Regular Quantiles}
\label{sec:quantiles}

\subsection{Mixture Model}
\label{sec:gmm}

\subsection{Comparison Metrics}
\label{sec:metrics}

[Lay out all metrics considered, pros and cons]

\subsubsection{RMS}
\label{sec:rms}

\subsubsection{KLD}
\label{sec:kld}

\subsubsection{Moments}
\label{sec:moments}

\subsubsection{KS}
\label{sec:ks}

% ----------------------------------------------------------------------

\section{Photo-z Test Data}
\label{sec:data}

\subsection{Simulation}
\label{sec:mock}

["Observed" $z_{s}$ vs. $z_{p}$ (from Buzzard)]

\begin{figure}
%\includegraphics[width=0.9\columnwidth]{}
\caption{[plot of $z_{s}$ vs. $z_{p}$]\label{fig:observed}}
\end{figure}

[Fit $z_{s}$ vs. $z_{p}$ with bivariate GMM]

\begin{figure}
%\includegraphics[width=0.9\columnwidth]{}
\caption{[$z_{s}$ vs. $z_{p}$ contour plot after GMM fit]\label{fig:fit}}
\end{figure}

[Specify true $n(z)$]

\begin{figure}
%\includegraphics[width=0.9\columnwidth]{}
\caption{[chosen true $n(z)$]\label{fig:nz}}
\end{figure}

\subsection{Catalog generation}
\label{sec:catalogs}

[Choose catalog size $N$ and sample $z_{s}$ values from true $n(z)$]

[Evaluate $N$ "true" $p(z)$s from horizontal cuts in fitted $z_{s}$ vs. $z_{p}$; these are \texttt{qp.catalog} objects comprised of \texttt{qp.PDF} objects defined by original GMM formulae]

\begin{figure}
%\includegraphics[width=0.9\columnwidth]{}
\caption{[Some example $p(z)$s]\label{fig:pzs}}
\end{figure}

[Convert into parametrizations as separate catalogs]

% ----------------------------------------------------------------------

\section{Science Metrics}
\label{sec:science}

\subsection{$n(z)$}
\label{sec:nz}

\begin{align}
\hat{n}(z) &= \frac{1}{N}\sum_{i=1}^{N}p_{i}(z)
\label{eq:nz}
\end{align}

\subsection{$\Sigma_{crit}^{-1}$}
\label{sec:sigma}

\citet{bonnett_redshift_2015}

\begin{align}
\Sigma_{crit}^{-1} &= \frac{4\pi G}{c}\frac{D_{ls}D_{l}}{D_{s}}
\label{eq:sigma}
\end{align}

\begin{align}
\left\langle\Sigma_{crit}^{-1}\right\rangle(z_{lens}) &= \sum_{i=1}^{N}p(z_{source, i})\Sigma_{crit}^{-1}(z_{lens}, z_{source, i})
\label{eq:meansigma}
\end{align}

% ----------------------------------------------------------------------

\section{Results}
\label{sec:results}

%\figref{example} shows an example figure, referred to with the \verb=\figref= command and the \code{fig:} prefix.
%
%\begin{figure}
%\includegraphics[width=0.9\columnwidth]{example.png}
%\caption{An example figure: the LSST DESC logo, copied from \code{.logos/desc-logo.png} into \code{figures/example.png}. \label{fig:example}}
%\end{figure}

[Multipanel key plot commentary]

\begin{figure}
%\includegraphics[width=0.9\columnwidth]{}
\caption{[number of stored floats $N_{f}$ vs. metric, one curve per approximation method]\label{fig:results}}
\end{figure}

% ----------------------------------------------------------------------

%\section{Discussion}
%\label{sec:discussion}
%
%If you are planning on committing your paper to GitHub, it's a good idea to write your tex as one sentence per line.
%This allows for an easier \code{diff} of changes.
%It also makes sense to think of latex as \emph{code}, and sentences as logical statements, occupying one line each.
%Each line must ``compile'' in the mind of the reader.


% ----------------------------------------------------------------------

\section{Conclusions \& Future Directions}
\label{sec:conclusions}

%Here's a summary of what we just reported.
%
%We can draw the following well-organized and neatly-formatted conclusions:
%\begin{itemize}
%  \item This is important.
%  \item We can measure some number with some precision.
%  \item This has some implications.
%\end{itemize}
%
%Here are some parting thoughts.


% ----------------------------------------------------------------------

\subsection*{Acknowledgments}

%Here is where you should add your specific acknowledgments, remembering that some standard thanks will be added via the \code{acknowledgments.tex} and \code{contributions.tex} files.

%
%This is the text imported from \code{acknowledgments.tex}, and will be replaced by some standard LSST DESC boilerplate at some point.
%


Author contributions are listed below. \\
A.I.~Malz: Initiated project, led development work. \\
P.J.~Marshall: Advised on statistics, and project design and management. \\
S.J.~Schmidt: Produced the PDFs for the fainter mock catalog. \\
M.L.~Graham: Produced the photometry and PDFs for the brighter mock catalog. \\
J.~DeRose: Produced the photometry for the fainter mock catalog. \\
R.~Wechsler: Supervised production of the fainter mock catalog. \\


%{\it Facilities:} \facility{LSST}

% Include both collaboration papers and external citations:
\bibliography{lsstdesc,main}

\end{document}
% ======================================================================
% 
